\documentclass[12pt]{article}
\usepackage[utf8]{inputenc}
\usepackage{amsmath, amssymb, amsthm}
\usepackage{geometry}
\geometry{a4paper, margin=1in}

\title{Mathematical Foundations: Pythagorean Triples}
\author{Study Notes}
\date{\today}

\begin{document}

\maketitle

\section{Introduction}
A \textbf{Pythagorean Triple} is a set of three positive integers $(a, b, c)$ satisfying the Diophantine equation:
\begin{equation}
    a^2 + b^2 = c^2
\end{equation}
These triples describe the integer side lengths of right-angled triangles.

\section{Classification}
\begin{itemize}
    \item \textbf{Primitive Pythagorean Triples (PPT):} A triple where $\gcd(a, b, c) = 1$. Every PPT can be used to generate infinitely many imprimitive triples by multiplying each term by a constant $k$.
    \item \textbf{Parity Properties:} In any PPT, the following must hold:
    \begin{itemize}
        \item One leg is always odd, and one leg is always even.
        \item The hypotenuse $c$ is always odd.
    \end{itemize}
\end{itemize}

\section{Euclid's Formula}
Euclid's formula is a fundamental method for generating all PPTs. For any two integers $m$ and $n$:
\begin{align*}
    a &= m^2 - n^2 \\
    b &= 2mn \\
    c &= m^2 + n^2
\end{align*}
To ensure the triple is \textbf{primitive}, the parameters must satisfy:
\begin{enumerate}
    \item $m > n > 0$
    \item $\gcd(m, n) = 1$
    \item $m+n$ is odd (one is even, one is odd)
\end{enumerate}

\section{Divisibility Rules}
Every Pythagorean triple $(a, b, c)$ obeys the following:
\begin{itemize}
    \item $3 \mid ab$ (At least one leg is divisible by 3)
    \item $4 \mid ab$ (At least one leg is divisible by 4)
    \item $5 \mid abc$ (At least one of the three sides is divisible by 5)
    \item The product $abc$ is always divisible by 60.
\end{itemize}

\section{Complex Number Derivation}
Pythagorean triples can be derived using \textbf{Gaussian Integers}. Consider a Gaussian integer $z = m + ni \in \mathbb{Z}[i]$. Squaring $z$ yields:
\begin{equation}
    z^2 = (m + ni)^2 = (m^2 - n^2) + (2mn)i
\end{equation}
The norm of $z^2$ is $|z^2| = |z|^2 = (m^2 + n^2)$. This mapping provides the algebraic foundation for Euclid's formula.

\section{Geometric Perspective}
Dividing $a^2 + b^2 = c^2$ by $c^2$ yields the unit circle equation:
\begin{equation}
    \left(\frac{a}{c}\right)^2 + \left(\frac{b}{c}\right)^2 = 1
\end{equation}
Thus, finding PPTs is equivalent to finding points on the unit circle $x^2 + y^2 = 1$ where $x$ and $y$ are both \textbf{rational numbers}.

\end{document}